\paragraph{E.4j (Refuted): integral all-twist purity with the verified finite-field groups and ordinary finite coefficients.}
There is a further obstruction at a finite residue prime, independent
of the omitted real sign in E.4b. The following requirements
cannot hold simultaneously. Let
\[
 \overline X=\overline{\operatorname{Spec}\mathbb Z},\qquad
 \overline U=\overline X-\{2\},\qquad z=\operatorname{Spec}\mathbb F_2,
\]
where the real place is retained, and suppose complexes $C_T^n$ have:
the finite-field values E.4a at $z$; the proper twist-zero complex
$C_{\overline X}^0$ of E.4b (or its sign-deleted version E.4f);
an integral purity/localization triangle
\[
 C_z^{-1}[-2]\longrightarrow C_{\overline X}^0
       \longrightarrow C_{\overline U}^0\longrightarrow C_z^{-1}[-1];
 \tag{E.4.10}
\]
and a finite-coefficient comparison
\[
 C_{\overline U}^0\otimes^L\mathbb Z/2
   \simeq R\Gamma(\overline U_{\mathrm{et}},\mathbb Z/2),
 \tag{E.4.11}
\]
with the constant sheaf on the Artin--Verdier site. Already their
degree-one consequences contradict one another.

\begin{proof}
For $q=2,n=-1$, E.4a says that every cohomology group of $C_z^{-1}$
is zero: its possible degree-one group has order $2^1-1=1$.
Thus (E.4.10) makes restriction a quasi-isomorphism.
But E.4b gives
\[
 C_{\overline X}^0\simeq
 \mathbb Z\oplus(\mathbb Z/2)[-3],
 \qquad
 H^1(C_{\overline X}^0\otimes^L\mathbb Z/2)=0.
\]
The sign-deleted choice $\mathbb Z$ has the same degree-one vanishing.
Hence (E.4.10)--(E.4.11) would imply
$H^1(\overline U_{\mathrm{et}},\mathbb Z/2)=0$.

There is an explicit nontrivial torsor contradicting that vanishing.
Put $R=\mathbb Z[1/2]$ and
\[
 B=R[t]/(t^2-2),\qquad \sigma(t)=-t.
\]
This is a finite free $R$-algebra of rank two, and it is \'etale
because $2t$ is a unit: $(2t)(t/4)=1$ in $B$.
The two factors $T-t$ and $T+t$ in $B[T]$ are comaximal, since
$2t$ is a unit. The Chinese remainder theorem therefore identifies
$B\otimes_R B$ with $B\times B$ by the two automorphisms
$1,\sigma$. Thus $\operatorname{Spec}B$ is a right torsor under the
constant group $\mathbb Z/2$ over $\operatorname{Spec}R$.
It has no section: a section would supply $a\in R\subset\mathbb Q$
with $a^2=2$, impossible because its $2$-adic valuation would satisfy
$2v_2(a)=1$. Equivalently $B$ has fraction field $\mathbb Q(\sqrt2)$
and is not the product algebra of a trivial torsor.

The extension is totally real and splits at the real completion:
$\mathbb Q(\sqrt2)\otimes\mathbb R\simeq\mathbb R\times\mathbb R$.
Attaching both real places therefore extends this torsor to the
Artin--Verdier site with the real point retained. It remains nontrivial,
since a section there would restrict to one over $R$.
Its class in $H^1(\overline U_{\mathrm{et}},\mathbb Z/2)$ is consequently
nonzero. This contradicts the required vanishing and proves the claim.
\end{proof}

The coefficient convention in (E.4.11) is concrete: in the proper
Flach--Morin theory, \citet[Proposition 6.11]{D_FlachMorin2018}
identifies $\mathbb Z(0)$ with the constant sheaf, and their
Lemma 3.13 identifies derived reduction with its finite-coefficient
\'etale cohomology. We have stated the analogous open comparison as
a requirement in order to prove its incompatibility with (E.4.10);
we have not misapplied their proper-scheme functoriality theorem to
the open immersion.

This counterexample explains a specific failed route toward the
all-twist integral extension. Geisser's Dedekind-ring purity theorem
has a truncation and a residue-characteristic qualification
\citep[Theorem 1.2(1)]{D_GeisserDedekind2004}. The missing qualification
cannot be removed in the stated coefficient theory. A corrected
integral localization theory must retain a nonzero contribution from
wild ramification at $2$, or alter (E.4.11), rather than use only the
finite-field complex $C_z^{-1}[-2]$ as its support term. Inverting $2$
removes this particular obstruction, and the rational result E.4g
is unaffected. The transcription 15:174 explicitly applies its displayed
purity sequence to $M=\mathbb Q(n)$, so this refutes an integral
extrapolation, not that rational assertion in the manuscript.

\paragraph{E.4k (Refuted): finite generation on opens together with the full coefficient comparison and rational restriction pattern.}
The obstruction persists without assuming the integral purity triangle.
There is no complex $C$ of abelian groups for which all three of the
following hold: $H^1(C)$ and $H^2(C)$ are finitely generated;
$H^1(C\otimes\mathbb Q)=H^2(C\otimes\mathbb Q)=0$;
and, for every $r\geq1$,
\[
 H^1(C\otimes^L\mathbb Z/2^r)
      \simeq H^1(\overline U_{\mathrm{et}},\mathbb Z/2^r),
 \qquad \overline U=
 \overline{\operatorname{Spec}\mathbb Z}-\{2\}.
 \tag{E.4.12}
\]
Thus the verified rational twist-zero restriction pattern cannot be
combined on this open with finite generation and every usual finite
coefficient comparison. No multiplicative or duality axiom is needed
for this refutation.

\begin{proof}
The first two assumptions make $A=H^1(C)$ and $B=H^2(C)$ finite.
The coefficient triangle supplies the exact sequence
\[
 0\longrightarrow A/2^rA
   \longrightarrow H^1(C\otimes^L\mathbb Z/2^r)
   \longrightarrow B[2^r]\longrightarrow0.
\]
Consequently the order of its middle group is at most $|A|\,|B|$,
independently of $r$.

We construct a class of exact order $2^r$ on the right of (E.4.12).
Set $m=r+2$ and let
\[
 K_m=\mathbb Q(\zeta_{2^m}+\zeta_{2^m}^{-1}),\qquad
 R=\mathbb Z[1/2],\qquad
 B_m=\bigl(R[T]/(T^{2^{m-1}}+1)\bigr)^{T\mapsto T^{-1}}.
\]
The polynomial is $\Phi_{2^m}(T)$; replacing $T$ by $T+1$ makes
it Eisenstein at $2$. The full cyclotomic algebra therefore has a
field as its generic fibre. It is finite \'etale over $R$, since
$T$ and $2^{m-1}T^{2^{m-1}-1}$ are units in the algebra.
The automorphisms $T\mapsto T^a$, for odd $a$ modulo $2^m$, give
its Galois torsor structure: on every geometric fibre the roots have
exact order $2^m$, and that group acts freely and transitively.
Taking the quotient by $a=-1$ shows that $\operatorname{Spec}B_m$
is a connected finite \'etale torsor under
\[
 G_m=(\mathbb Z/2^m\mathbb Z)^\times/\{\pm1\}
       \simeq\mathbb Z/2^{m-2}\mathbb Z.
\]
For the displayed cyclic identification one may use $5$ as a generator.
Indeed induction by squaring gives
$v_2(5^{2^s}-1)=s+2$ for every $s\geq0$, so $5$ has exact order
$2^{m-2}$ modulo $2^m$ and generates the units congruent to $1$
modulo $4$. Every odd residue is uniquely a sign times such a unit.

The generic field of $B_m$ is $K_m$, which is totally real.
Attaching all of its real places gives a connected torsor on
$\overline U_{\mathrm{et}}$ under $G_m\simeq\mathbb Z/2^r$.
Its monodromy character onto that cyclic group is surjective,
so its class in $H^1(\overline U_{\mathrm{et}},\mathbb Z/2^r)$
has exact order $2^r$. In particular the order of this cohomology
group is at least $2^r$, contradicting the preceding uniform bound
as $r$ increases. This proves the assertion.
\end{proof}

A possible integral extension retaining the rational pattern and all
these finite comparisons must therefore allow non-finitely-generated
torsion in the adjacent integral groups, or change one of those
requirements. For instance, a divisible $2$-primary torsion summand
can have nonzero $2^r$-torsion while having zero rationalization;
finite generation had excluded exactly this possibility.
The proposition does not assert nonexistence of all integral open
cohomology theories. It identifies the precise additional condition
that fails, beyond the already refuted exact proper sign table.
